\documentclass[11pt,letterpaper]{letter}
\addtolength{\oddsidemargin}{-.5in}%
\addtolength{\evensidemargin}{-.5in}%
\addtolength{\textwidth}{1.3in}%
\addtolength{\textheight}{1.5in}%
\addtolength{\topmargin}{-1in}%
\address{José Á. Sánchez Gómez\\
%Assistant Professor\\
Department of Statistics\\
University of California Riverside\\
Riverside, CA 92521\\
josesa@ucr.edu }

\signature{José Á. Sánchez Gómez}

\begin{document}
%\date{\today}
\date{}


\begin{letter}{%Professors A. Delaigle and S. Wood \\
Editors\\
Journal of the Americal Statistical Association}

\opening{Dear Editors,}

\bigskip
re: ``Quantifying Moderator Effects on High-Dimensional Multi-Subject Time Series'' by José Á. Sánchez Gómez, Holly O'Rourke and Gene Brewer. 

\bigskip

%Thank you very much for the invitation to the contribution to Wiley Handbook of Statistical Computing. 
Could you please consider the above article for possible publication in {\em JASA}? 

	Researchers across many areas use time-series data to estimate relationships among variables over time. These relationships can be modeled via Granger-causal networks (GCNs), where nodes correspond to variables and directed edges indicate lagged temporal relationships among them. When time-series data are gathered over multiple subjects, GCN models must account for the presence of shared and subject-specific edges across the GCNs. While existing methods can estimate GCNs on multi-subject data, they often lack interpretable explanations for the source of the observed heterogeneity. In the present paper, we propose the moderator vector autoregressive model (MOD-VAR) for interpretable estimation of GCNs in multi-subject time-series data. To this end, we decompose the GCNs into shared, moderator-dependent, and subject-specific components. The use of moderator-dependent components can enhance both the quantification and interpretability of heterogeneous GCNs. Our theoretical results establish that MOD-VAR improves multi-subject GCN estimation in high-dimensional regimes with substantial heterogeneity. Furthermore, extensive numerical simulations demonstrate that MOD-VAR outperforms existing methods in terms of both GCN estimation and forecasting in a variety of settings.  

Thank you very much for your consideration. I look forward to hearing from you.

\closing{Yours sincerely}

\end{letter}
\end{document}
